\section{Capacity feasibility and the crossing penalty}
\subsection*{Bin Packing inside the routing problem}
Bin Packing decision gives item sizes $a_1,\ldots,a_n$, bin capacity $C$ and at most $k$ bins.
Turn each item into a clinic with demand $q_i=a_i$, and set $Q=C$ and $K=k$. Use distance 1
between distinct locations, since distance is irrelevant to this feasibility question. If an
item exceeds $C$, the input is immediately a NO instance.

A packing gives one route per bin, visiting its items in any order. Loads are at most $Q$ and
there are at most $K$ routes. Conversely, a feasible route assignment gives bins containing
those same demands. Both directions preserve the answer, and the construction is polynomial.
A route assignment is also easy to check. Thus even the distance-free fleet-feasibility problem
is NP-complete, by reduction from Bin Packing [1].

For example, demands $(6,6,6)$ total 18. Two drones of capacity 10 have total capacity 20, but
cannot serve these demands: no two fit together. Total capacity alone does not guarantee feasibility.

This matters in the code. CW respects individual capacities but may produce too many routes.
GA checks decoded route counts as well as loads. TS may temporarily overload routes, but saves
a final answer only when every constraint holds.

\subsection*{Why a fixed crossing cost preserves the metric}
Let $e(i,j)$ be unrounded Euclidean distance and $c(i,j)$ be 1 when the endpoints are on different
sides of the Bosphorus, and 0 otherwise. Define
\[
 d(i,j)=e(i,j)+P\,c(i,j),\qquad P\ge0.
\]
Symmetry and zero self-distance are preserved. For the triangle inequality, consider going from
$i$ to $k$ via $j$. If $i$ and $k$ are on different sides, $j$ must be on one side, so at least
one leg of the detour also crosses. If they are on the same side, their direct crossing cost is
zero. In either case,
\[
 c(i,k)\le c(i,j)+c(j,k).
\]
Adding this inequality, multiplied by $P$, to the Euclidean triangle inequality gives
$d(i,k)\le d(i,j)+d(j,k)$. The crossing penalty therefore preserves the metric property.

The synthetic boundary is $x=500$. The depot at $(500,500)$ belongs to Asia because the rule is
$x\ge500$. A route serving Europe must cross and return.

\subsection*{Benchmark rounding is a separate convention}
CVRPLIB requires Euclidean distances rounded to the nearest integer for comparison with its
reference costs. Rounding can introduce small triangle-inequality violations. The exact proof
above applies to unrounded MediDrone costs. The benchmark evaluator follows the benchmark's
integer convention.

\takeaway{A plan must pass both capacity and fleet checks. The crossing proof uses a simple
fact: a direct trip cannot require more side changes than a detour.}
